\chapter{A Quarter of an Hour's Delay}

\lettrine{F}{ouquet}, on leaving his house for the second time that day, felt himself less heavy and less disturbed than might have been expected. He turned towards Pélisson, who was meditating in the corner of the carriage some good arguments against the violent proceedings of Colbert.

<My dear Pélisson,> said Fouquet, <it is a great pity you are not a woman.>

<I think, on the contrary, it is very fortunate,> replied Pélisson, <for, monseigneur, I am excessively ugly.>

<Pélisson! Pélisson!> said the superintendent, laughing: <You repeat too often, you are <ugly>, not to leave people to believe that it gives you much pain.>

<In fact it does, monseigneur, much pain; there is no man more unfortunate than I: I was handsome, the small-pox rendered me hideous; I am deprived of a great means of attraction; now, I am your principal clerk, or something of that sort; I take great interest in your affairs, and if, at this moment, I were a pretty woman, I could render you an important service.>

<What?>

<I would go and find the concierge of the Palais. I would seduce him, for he is a gallant man, extravagantly partial to women; then I would get away our two prisoners.>

<I hope to be able to do so myself, although I am not a pretty woman,> replied Fouquet.

<Granted, monseigneur; but you are compromising yourself very much.>

<Oh!> cried Fouquet, suddenly, with one of those secret transports which the generous blood of youth, or the remembrance of some sweet emotion, infuses into the heart. <Oh! I know a woman who will enact the personage we stand in need of, with the lieutenant-governor of the concierge.>

<And, on my part, I know fifty, monseigneur; fifty trumpets, which will inform the universe of your generosity, of your devotion to your friends, and, consequently, will ruin you sooner or later in ruining themselves.>

<I do not speak of such women, Pélisson; I speak of a noble and beautiful creature who joins to the intelligence and wit of her sex the valour and coolness of ours; I speak of a woman, handsome enough to make the walls of a prison bow down to salute her, discreet enough to let no one suspect by whom she has been sent.>

<A treasure!> said Pélisson; <you would make a famous present to monsieur the governor of the concierge! \swear{Peste!} monseigneur, he might have his head cut off; but he would, before dying, have had such happiness as no man had enjoyed before him.>

<And I add,> said Fouquet, <that the concierge of the Palais would not have his head cut off, for he would receive of me my horses, to effect his escape, and five hundred thousand livres wherewith to live comfortably in England: I add, that this lady, my friend, would give him nothing but the horses and the money. Let us go and seek her, Pélisson.>

The superintendent reached forth his hand towards the golden and silken cord placed in the interior of his carriage, but Pélisson stopped him. <Monseigneur,> said he, <you are going to lose as much time in seeking this lady as Columbus took to discover the new world. Now, we have but two hours in which we can possibly succeed; the concierge once gone to bed, how shall we get at him without making a disturbance? When daylight dawns, how can we conceal our proceedings? Go, go yourself, monseigneur, and do not seek either woman or angel to-night.>

<But, my dear Pélisson, here we are before her door.>

<What! before the angel's door?>

<Why, yes.>

<This is the hotel of Madame de Bellière!>

<Hush!>

<Ah! Good Lord!> exclaimed Pélisson.

<What have you to say against her?>

<Nothing, alas! and it is that which causes my despair. Nothing, absolutely nothing. Why can I not, on the contrary, say ill enough of her to prevent your going to her?>

But Fouquet had already given orders to stop, and the carriage was motionless. <Prevent me!> cried Fouquet; <why, no power on earth should prevent my going to pay my compliments to Madame de Plessis-Bellière; besides, who knows that we shall not stand in need of her!>

<No, monseigneur, no!>

<But I do not wish you to wait for me, Pélisson,> replied Fouquet, sincerely courteous.

<The more reason I should, monseigneur; knowing that you are keeping me waiting, you will, perhaps, stay a shorter time. Take care! You see there is a carriage in the courtyard: she has some one with her.> Fouquet leaned towards the steps of the carriage. <One word more,> cried Pélisson; <do not go to this lady till you have been to the concierge, for Heaven's sake!>

<Eh! five minutes, Pélisson,> replied Fouquet, alighting at the steps of the hotel, leaving Pélisson in the carriage, in a very ill-humour. Fouquet ran upstairs, told his name to the footman, which excited an eagerness and a respect that showed the habit the mistress of the house had of honouring that name in her family. <Monsieur le surintendant,> cried the marquise, advancing, very pale, to meet him; <what an honour! what an unexpected pleasure!> said she. Then, in a low voice, <Take care!> added the marquise, <Marguerite Vanel is here!>

<Madame,> replied Fouquet, rather agitated, <I came on business. One single word, and quickly, if you please!> And he entered the salon. Madame Vanel had risen, paler, more livid, than Envy herself. Fouquet in vain addressed her, with the most agreeable, most pacific salutation; she only replied by a terrible glance darted at the marquise and Fouquet. This keen glance of a jealous woman is a stiletto which pierces every cuirass; Marguerite Vanel plunged it straight into the hearts of the two confidants. She made a courtesy to her friend, a more profound one to Fouquet, and took leave, under pretence of having a number of visits to make, without the marquise trying to prevent her, or Fouquet, a prey to anxiety, thinking further about her. She was scarcely out of the room, and Fouquet left alone with the marquise, before he threw himself on his knees, without saying a word. <I expected you,> said the marquise, with a tender sigh.

<Oh! no,> cried he, <or you would have sent away that woman.>

<She has been here little more than half an hour, and I had no expectation she would come this evening.>

<You love me just a little, then, marquise?>

<That is not the question now; it is of your danger; how are your affairs going on?>

<I am going this evening to get my friends out of the prisons of the Palais.>

<How will you do that?>

<By buying and bribing the governor.>

<He is a friend of mine; can I assist you, without injuring you?>

<Oh! marquise, it would be a signal service; but how can you be employed without your being compromised? Now, never shall my life, my power, or even my liberty, be purchased at the expense of a single tear from your eyes, or of one frown of pain upon your brow.>

<Monseigneur, no more such words, they bewilder me; I have been culpable in trying to serve you, without calculating the extent of what I was doing. I love you in reality, as a tender friend; and as a friend, I am grateful for your delicate attentions—but, alas!—alas! you will never find a mistress in me.>

<Marquise!> cried Fouquet, in a tone of despair; <why not?>

<Because you are too much beloved,> said the young woman, in a low voice; <because you are too much beloved by too many people—because the splendour of glory and fortune wound my eyes, whilst the darkness of sorrow attracts them; because, in short, I, who have repulsed you in your proud magnificence; I who scarcely looked at you in your splendour, I came, like a mad woman, to throw myself, as it were, into your arms, when I saw a misfortune hovering over your head. You understand me now, monseigneur? Become happy again, that I may remain chaste in heart and in thought: your misfortune entails my ruin.>

<Oh! madame,> said Fouquet, with an emotion he had never before felt; <were I to fall to the lowest degree of human misery, and hear from your mouth that word which you now refuse me, that day, madame, you will be mistaken in your noble egotism; that day you will fancy you are consoling the most unfortunate of men, and you will have said, I love you, to the most illustrious, the most delighted, the most triumphant of the happy beings of this world.>

He was still at her feet, kissing her hand, when Pélisson entered precipitately, crying, in very ill-humour, <Monseigneur! madame! for Heaven's sake! excuse me. Monseigneur, you have been here half an hour. Oh! do not both look at me so reproachfully. Madame, pray who is that lady who left your house soon after monseigneur came in?>

<Madame Vanel,> said Fouquet.

<Ha!> cried Pélisson, <I was sure of that.>

<Well! what then?>

<Why, she got into her carriage, looking deadly pale.>

<What consequence is that to me?>

<Yes, but what she said to her coachman is of consequence to you.>

<Kind heaven!> cried the marquise, <what was that?>

<To M.~Colbert's!> said Pélisson, in a hoarse voice.

<\swear{Bon Dieu!}—begone, begone, monseigneur!> replied the marquise, pushing Fouquet out of the salon, whilst Pélisson dragged him by the hand.

<Am I, then, indeed,> said the superintendent, <become a child, to be frightened by a shadow?>

<You are a giant,> said the marquise, <whom a viper is trying to bite in the heel.>

Pélisson continued to drag Fouquet to the carriage. <To the Palais at full speed!> cried Pélisson to the coachman. The horses set off like lightening; no obstacle relaxed their pace for an instant. Only, at the arcade Saint-Jean, as they were coming out upon the Place de Grève, a long file of horsemen, barring the narrow passage, stopped the carriage of the superintendent. There was no means of forcing this barrier; it was necessary to wait till the mounted archers of the watch, for it was they who stopped the way, had passed with the heavy carriage they were escorting, and which ascended rapidly towards the Place Baudoyer. Fouquet and Pélisson took no further account of this circumstance beyond deploring the minute's delay they had thus to submit to. They entered the habitation of the concierge du Palais five minutes after. That officer was still walking about in the front court. At the name of Fouquet, whispered in his ear by Pélisson, the governor eagerly approached the carriage, and, hat in hand, was profuse in his attentions. <What an honour for me, monseigneur,> said he.

<One word, monsieur le governeur, will you take the trouble to get into my carriage?> The officer placed himself opposite Fouquet in the coach.

<Monsieur,> said Fouquet, <I have a service to ask of you.>

<Speak, monseigneur.>

<A service that will be compromising for you, monsieur, but which will assure to you forever my protection and my friendship.>

<Were it to cast myself into the fire for you, monseigneur, I would do it.>

<That is well,> said Fouquet; <what I require is much more simple.>

<That being so, monseigneur, what is it?>

<To conduct me to the chamber of Messieurs Lyodot and D'Eymeris.>

<Will monseigneur have the kindness to say for what purpose?>

<I will tell you that in their presence, monsieur; at the same time that I will give you ample means of palliating this escape.>

<Escape! Why, then, monseigneur does not know?>

<What?>

<That Messieurs Lyodot and D'Eymeris are no longer here.>

<Since when?> cried Fouquet, in great agitation.

<About a quarter of an hour.>

<Whither have they gone, then?>

<To Vincennes—to the donjon.>

<Who took them from here?>

<An order from the king.>

<Oh! woe! woe!> exclaimed Fouquet, striking his forehead. <Woe!> and without saying a single word more to the governor, he threw himself back into his carriage, despair in his heart, and death on his countenance.

<Well!> said Pélisson, with great anxiety.

<Our friends are lost. Colbert is conveying them to the donjon. They crossed our path under the arcade Saint-Jean.>

Pélisson, struck as by a thunderbolt, made no reply. With a single reproach he would have killed his master. <Where is monseigneur going?> said the footman.

<Home—to Paris. You, Pélisson, return to Saint-Mande, and bring the Abbé Fouquet to me within an hour. Begone!> 